\begin{frame}{``TD의 분산이 작다''는 무엇이냐}
  \begin{itemize}
    \item 표의 ``분산 큼/작음''은 \textbf{갱신 한 번에 쓰이는 목표값이
      얼마나 흔들리느냐}의 이야기.
    \item 초기 $V(s')=0$이라 TD 목표값이 1로 시작해 범위가 1$\sim$11.7까지
      벌어지지만, 표준편차는 MC의 절반 이하.
    \item 편향 쪽도 숫자로: 첫 에피소드에서 상태 4의 목표값
      $1+V(5)=1$은 참값 8과 \textbf{크게} 어긋남.
  \end{itemize}
  \vskip0.6em
  \begin{block}{부트스트래핑의 가격표}
    $V(s')$가 부정확한 한, TD의 목표값은 참값을 \textbf{체계적으로
    어긋나 있는}(= 편향된) 상태다. MC의 목표값 $G$는 어떤 에피소드가
    아무리 운이 없어도 \textbf{평균}으로는 정확하다(불편).
    ``작은 분산''을 얻는 대가로 ``임시적인 편향''을 쓴다 ---
    그리고 그 편향은 $V$가 수렴할수록 \textbf{자동으로 사라진다}.
  \end{block}
  \vskip0.4em
  {\footnotesize 최종 추정치 전체의 산포는 $\alpha$ 스케줄과 샘플 수의 문제
  (아래 ``자주 하는 실수'' 세 번째에서 계속).}
\end{frame}
